\documentclass[11pt]{article}
\usepackage[a4paper,margin=27mm]{geometry}
\usepackage{amsmath,amssymb,amsthm}
\usepackage{microtype}
\usepackage{booktabs,tabularx}
\usepackage[hidelinks]{hyperref}
\newtheorem{definition}{Definition}[section]
\newtheorem{proposition}[definition]{Proposition}
\newtheorem{theorem}[definition]{Theorem}
\newtheorem{corollary}[definition]{Corollary}
\newtheorem{remark}[definition]{Remark}
\newcommand{\Wcal}{\mathcal W}
\newcommand{\Ccal}{\mathfrak C}
\newcommand{\Hminus}{\mathcal H_-}
\newcommand{\Pcal}{\mathfrak H_-}

\hypersetup{
  pdftitle={WR-IV v0.4: Response-Conditioned Global Completion and the Status of the Past},
  pdfauthor={A. A. Malachevsky},
  pdfsubject={World Realizability and Epistemic Horizons; review draft},
  pdfkeywords={history projection, realizability, structural walls, Brownian bridge, smoothing}
}

\title{\textbf{WR-IV Mathematical Spine v0.4}\\[4pt]\large Response-Conditioned Global Completion and the Status of the Past}
\author{A. A. Malachevsky\\ORCID: 0009-0008-6009-3196\\World Realizability \& Epistemic Horizons (WREH)}
\date{8 October 2026}
\begin{document}
\maketitle

\begin{abstract}
WR-IV studies historical ambiguity inside response-conditioned global completions. Given present data $D$, the globally compatible world class $\Ccal(D)$ is projected onto declared past-history spaces. This separates world-level nonuniqueness from past-level nonuniqueness and from nonuniqueness of observable historical records. Accumulated hard constraints give nested completion classes, possibly equal, without altering retained histories. The WR-II realizability gate distinguishes an empty historical class from an ambiguous one. Transfer of WR-III structural walls to historical families requires a topology-preserving projection of the response-indexed family. A natural-extension benchmark for the doubling map shows that an exactly known present state may have a unique deterministic future while admitting a Cantor family of compatible full pasts. A second, physical inverse-problem benchmark uses fixed-interval smoothing of a diffusion: an endpoint observation reduces interior-state variance while the conditioned path law remains non-degenerate. Full-history compatibility and uncertainty about a fixed earlier fragment are distinguished. No claim of physical past creation, retrocausality, or branching ontology is made.
\end{abstract}

\section{Claim ceiling}
WR-IV inherits response equivalence and completion-fibre language from
WR-I--III \cite{WRI,WRII,WRIII}. The realizability gate below is an
inherited existence discipline, and the structural-wall bridge is a
conditional transfer of WR-III geometry, not a new existence or
singularity theory. It assumes a model in which a temporal
decomposition is already meaningful. It does not infer physical time from WREH itself. It does not claim observation creates or rewrites the past, that mathematical completion implies physical existence, or that noninvertibility implies a branching ontology.

\section{Response-conditioned completions}
Let
\[
\Ccal(D)=\{W\in\Wcal: W\text{ is globally compatible with }D\}.
\]
A declared historical projection is a map
\[
\pi_-:\Wcal\to\Hminus.
\]
\begin{definition}[Past-completion set]
\[
\Pcal(D):=\pi_-(\Ccal(D)).
\]
\end{definition}
\begin{definition}[Past rigidity]
$D$ is past-rigid if $\Pcal(D)$ is a singleton, and world-rigid if $\Ccal(D)$ is a singleton.
\end{definition}

\begin{proposition}[World rigidity implies past rigidity]
\label{prop:world-past}
If $D$ is world-rigid, then $D$ is past-rigid.
\end{proposition}
\begin{proof}
The image of a singleton under any map is a singleton.
\end{proof}

\begin{proposition}[The converse fails]
\label{prop:converse}
Past rigidity need not imply world rigidity.
\end{proposition}
\begin{proof}
Take $\Wcal=\{(h_0,0),(h_0,1)\}$, set $\Ccal(D)=\Wcal$,
and project both worlds onto $h_0$. Then $\Pcal(D)=\{h_0\}$,
although $\Ccal(D)$ has two elements.
\end{proof}


\subsection{The WR-II realizability gate}

For a declared protocol family $\mathcal P$, let
$R:\Wcal\to\mathcal Y=\prod_{P\in\mathcal P}Y_P$ be the full
response map. Use the profile representation of WR-II
\cite[Sections 3--7]{WRII}. For finite $\mathcal A\subseteq\mathcal P$,
write $R_{\mathcal A}=(R_P)_{P\in\mathcal A}$ and let
$y_{\mathcal A}$ denote the coordinate restriction of $y$. Then
\[
\begin{aligned}
\mathfrak C_{\mathrm{fin}}
&=\{y\in\mathcal Y:y_{\mathcal A}\in R_{\mathcal A}(\Wcal)
  \text{ for every finite }\mathcal A\subseteq\mathcal P\},\\
\mathfrak D_{\mathrm{fin}}
&=\mathfrak C_{\mathrm{fin}}\setminus R(\Wcal).
\end{aligned}
\]
For the full-profile data $D_y$, compatibility means
$\Ccal(D_y)=R^{-1}(y)$.

\begin{proposition}[Realizability precedes historical ambiguity]
\label{prop:realizability-gate}
For the total map $\pi_-:\Wcal\to\Hminus$,
\[
\Pcal(D_y)=\varnothing
\quad\Longleftrightarrow\quad
\Ccal(D_y)=\varnothing
\quad\Longleftrightarrow\quad
y\notin R(\Wcal).
\]
In particular, a full profile $y\in\mathfrak D_{\mathrm{fin}}$
has no compatible past in the declared model, despite realization
of every finite restriction.
\end{proposition}
\begin{proof}
An image under a total map is empty exactly when its domain is
empty. The fibre $R^{-1}(y)$ is empty exactly when $y$ is outside
the response image. The final statement uses the definition of
$\mathfrak D_{\mathrm{fin}}$.
\end{proof}

\begin{remark}[Empty class versus ambiguous past]
\label{rem:empty-past}
Emptiness is model-relative nonrealizability, not past rigidity,
physical disappearance of the past, or a failure of the Universe to
exist. A defect of a full profile does not make its finite-data fibres
empty. In the finite-support binary-world example of WR-II
\cite[Section 7]{WRII}, $y=(1,1,\ldots)$ has an empty full fibre,
but every finite restriction has a nonempty fibre. This is distinct
from conditioning on a probability-zero event with a nonempty fibre,
as in the Brownian endpoint example below.
\end{remark}

\section{Evidence refinement and non-rewriting}
Throughout refinement, $\Wcal$, the reference temporal cut, and $\pi_-$
are fixed. Write $D\preceq D'$ when $D'$ contains all hard constraints
of $D$ and possibly further constraints. This order does not include a
change of model or replacement of previous data. Refinement may leave
the completion set unchanged or make it empty; consistency is stated
separately when needed.
\begin{theorem}[Completion refinement]
\label{thm:completion-refinement}
If $D\preceq D'$, then
\[
\Ccal(D')\subseteq\Ccal(D)
\qquad\text{and}\qquad
\Pcal(D')\subseteq\Pcal(D).
\]
\end{theorem}
\begin{proof}
Every world satisfying the stronger data set also satisfies the weaker one; applying $\pi_-$ preserves inclusion.
\end{proof}
\begin{corollary}[Past rigidity persists under consistent refinement]
\label{cor:rigidity-persistence}
If $D$ is past-rigid, $D\preceq D'$, and $\Ccal(D')\neq\varnothing$, then $D'$ is past-rigid with the same past projection.
\end{corollary}
\begin{remark}[Non-rewriting principle]
Set refinement removes incompatible completions; it does not alter the histories carried by completions that remain. No physical rewriting of the past is implied.
\end{remark}

\section{Observable past versus hidden past}
Let $Q_-:\Hminus\to\mathcal O_-$ encode accessible historical records and define
\[
\mathfrak O_-(D)=Q_-(\Pcal(D)).
\]
\begin{proposition}[Observable rigidity criterion]
\label{prop:observable-rigidity}
$\mathfrak O_-(D)$ is a singleton if and only if $\Pcal(D)$ is
nonempty and $Q_-$ is constant on $\Pcal(D)$. Observable-past rigidity
can coexist with hidden-past ambiguity.
\end{proposition}
\begin{proof}
The criterion follows directly from the image definition. For the
counterexample, take $\Pcal(D)=\{h_1,h_2\}$ with $h_1\ne h_2$ and
$Q_-(h_1)=Q_-(h_2)$.
\end{proof}
Equality of records for one pair does not establish rigidity of a
larger past-completion set.


\subsection{When WR-III structural walls survive historical projection}

In a topological response model, let $R:\Wcal\to E=R(\Wcal)$
and $\pi_-:\Wcal\to\Hminus$ be continuous, with the declared
topologies. For an open $U\subseteq E$, give
\[
\mathcal H_U
=\{(y,h)\in U\times\Hminus:h\in\pi_-(R^{-1}(y))\}
\]
the subspace topology and let $p(y,h)=y$.
WR-III defines a structural wall by failure of local fibre-bundle
triviality over the attainable response image
\cite[Section 8]{WRIII}.

\begin{proposition}[Conditional transfer of structural walls]
\label{prop:wall-transfer}
If
\[
F:R^{-1}(U)\longrightarrow\mathcal H_U,
\qquad F(W)=(R(W),\pi_-(W))
\]
is a homeomorphism, then $R:R^{-1}(U)\to U$ and
$p:\mathcal H_U\to U$ have the same local-triviality locus.
Thus their structural walls coincide within $U$.
\end{proposition}
\begin{proof}
The identity $p\circ F=R$ makes $F$ an isomorphism of the two
families over $U$. Compose a local bundle trivialization with $F$
or $F^{-1}$ to transfer it in either direction.
\end{proof}

\begin{remark}[Projection can hide a structural wall]
\label{rem:wall-erasure}
The hypothesis concerns the whole response-indexed family, not
only a separate homeomorphism on each fibre. It is sufficient and
is not claimed necessary. Without such a bridge, a response wall
need not be a past wall. In the sphere-height example of WR-III
\cite[Section 10]{WRIII}, $R:S^2\to[-1,1]$, $R(x,y,z)=z$,
has circle fibres in the interior and singleton fibres at the
endpoints. Both endpoints are structural wall points. If the
declared historical projection is constant, $\pi_-(W)=h_0$, then
$\mathcal H_E=E\times\{h_0\}$ is a trivial historical family
with no wall. The example tests information loss under projection;
it supplies no physical temporal interpretation. Even an inherited
failure of local triviality does not by definition require a change
in the number of connected components. A response map is not a
physical forward-evolution map unless a separate bridge says so.
\end{remark}

\section{Deterministic future with ambiguous past}
Let $T:X\to X$ be deterministic. A full backward history ending at $x_0$ is
\[
(x_0,x_{-1},x_{-2},\ldots),
\qquad T(x_{-n})=x_{-(n-1)}.
\]
Define
\[
\widehat X=\{(x_0,x_{-1},\ldots)\in X^{\mathbb N_0}:
T(x_{-n})=x_{-(n-1)},\ n\ge1\},
\qquad p_0(\widehat x)=x_0.
\]
Histories over a fixed $x_0$ form the fibre $p_0^{-1}(x_0)$,
not the entire history space. For a continuous map on a topological
space, use the subspace topology inherited from the product.
These inverse-limit constructions are standard; the doubling-map
case is the dyadic solenoid \cite[Section 1]{ClarkFokkink2004}.
\begin{proposition}[Forward uniqueness and backward ambiguity]
\label{prop:forward-backward}
The forward orbit of $x_0$ is unique. Noninjectivity may allow
multiple full backward histories, but does not guarantee their
existence or multiplicity over every state.
\end{proposition}
\begin{proof}
The forward orbit is $(x_0,T(x_0),T^2(x_0),\ldots)$.
The doubling map below supplies an example of multiple full histories.
In contrast, for $T:\mathbb N_0\to\mathbb N_0$ with $T(0)=1$ and
$T(n)=n$ for $n\ge1$, no full history ends at $0$. Over $1$ the only
full history is $(1,1,\ldots)$, since a predecessor $0$ cannot be
continued backward.
\end{proof}

\section{Benchmark: doubling map}
Let $X=\mathbb R/\mathbb Z$ and $T(x)=2x\pmod1$. Each present state has two one-step predecessors.
\begin{theorem}[Cantor family of compatible full pasts]
\label{thm:cantor-pasts}
For every $x_0\in X$, the natural-extension fibre over $x_0$ is homeomorphic to $\{0,1\}^{\mathbb N}$.
\end{theorem}
\begin{proof}
Fix the representative $r_0\in[0,1)$ of $x_0$. Given
$\epsilon\in\{0,1\}^{\mathbb N}$, put
\[
r_n=(r_{n-1}+\epsilon_n)/2,\qquad x_{-n}=[r_n].
\]
Every history has unique representatives $r_n\in[0,1)$ and recovers
the bits by $\epsilon_n=2r_n-r_{n-1}\in\{0,1\}$. The coding is
therefore bijective. Each coordinate depends only on a finite prefix
of bits, so the coding is continuous. Its domain is compact and the
fibre is Hausdorff, hence the inverse is continuous. The identification
uses the chosen branch labels; no global continuous inverse branch on
the circle is asserted.
\end{proof}
\begin{corollary}[Finite backward records]
\label{cor:finite-backward}
Fixing any finite number of backward branch choices still leaves
a Cantor family of compatible full histories.
\end{corollary}
\begin{proof}
Deleting the finitely many fixed bit coordinates leaves a countably
infinite set of free coordinates, whose product is homeomorphic to
$\{0,1\}^{\mathbb N}$. For a fixed initial block these free coordinates
describe older histories.
\end{proof}

\section{Physical inverse benchmark: retrospective smoothing of a diffusion}

The doubling-map benchmark is exact and topological but not yet a physical inverse problem.
We now add a standard stochastic-state-estimation benchmark. Fixed-interval smoothing
estimates an earlier state using observations that occur after that state; in the linear
Gaussian case the classical Rauch--Tung--Striebel smoother gives a closed-form solution
\cite[Section 3.2]{RauchTungStriebel1965}
\cite[Section 12.2, Theorem 12.2]{SarkkaSvensson2023}.
The formulas below use direct Gaussian conditioning, not an
implementation or a new derivation of the RTS recursion.

\subsection{Gaussian retrospective contraction}

Fix a previously observed data value $A=a$ and work throughout
under its declared conditional law. Let $X\in\mathbb R^n$ denote a
past state and $Z\in\mathbb R^m$ later data. Under this fixed law,
assume $(X,Z)$ is jointly Gaussian with means $m_X,m_Z$ and covariance
\[
\begin{pmatrix}
P & C\\
C^\top & S
\end{pmatrix},
\qquad
S\succ0.
\]

\begin{theorem}[Later Gaussian data cannot increase past covariance]
\label{thm:gaussian-contraction}
The conditional covariance of the earlier state after incorporating the later data is
\[
P_{\mathrm{smooth}}
=
P-CS^{-1}C^\top.
\]
Hence
\[
0\preceq P_{\mathrm{smooth}}\preceq P.
\]
Moreover, the reduction is strict in every direction coupled to the later data through \(C\).
\end{theorem}

\begin{proof}
Set $E=X-m_X-CS^{-1}(Z-m_Z)$. Then
$\operatorname{Cov}(E,Z)=0$, so Gaussianity makes $E$ and $Z$
independent. The continuous Gaussian conditional kernel is
\[
X\mid Z=z\sim\mathcal N
\bigl(m_X+CS^{-1}(z-m_Z),\,P-CS^{-1}C^\top\bigr),
\]
where degenerate Gaussian laws are allowed. This is the standard
conditioning formula \cite[Appendix A.1, Lemma A.3]{SarkkaSvensson2023}.
In particular,
$P-CS^{-1}C^\top=\operatorname{Cov}(E)\succeq0$.
Since \(S^{-1}\succ0\),
\[
CS^{-1}C^\top\succeq0,
\]
so \(P_{\mathrm{smooth}}\preceq P\). For a vector \(v\),
\[
v^\top(P-P_{\mathrm{smooth}})v
=
(C^\top v)^\top S^{-1}(C^\top v),
\]
which is positive exactly when \(C^\top v\neq0\).
\end{proof}

\begin{definition}[Retrospective covariance gain]
Define
\[
\Delta P:=P-P_{\mathrm{smooth}}=CS^{-1}C^\top\succeq0.
\]
\end{definition}

\begin{remark}[Probabilistic extension]
All covariances above refer to the same fixed earlier data value
$A=a$. The theorem controls covariance, and does not assert nesting
of supports or of prior and posterior credible regions: their means
may differ. In the scalar nondegenerate Gaussian benchmark, central
intervals of a fixed probability become shorter, but may move outside
the prior interval. A zero coupling $C=0$ gives no reduction.
Positive residual covariance is not guaranteed by the general theorem;
it is established separately below. Outside Gaussian models,
pointwise monotonicity of conditional variance is not asserted.
Exact hard-constraint refinement and probabilistic smoothing remain
distinct layers.
\end{remark}

\subsection{Brownian bridge as an exactly solvable physical model}

Fix $x_0\in\mathbb R$, $T>0$, and diffusivity $\kappa>0$. Let
\[
X_t=x_0+\sqrt{2\kappa}\,B_t,\qquad 0\le t\le T.
\]
The coefficient $\kappa$ has units of length squared per unit time
and is distinct from the evidence symbol $D$.
Here $B$ is standard Brownian motion. This is an idealized physical
diffusion model, rather than an experimental validation of any specific
system; see \cite{KaratzasShreve1998} for Brownian-motion background.

With the deterministic initial state $X_0=x_0$ fixed, the marginal
law before the endpoint observation (the prior for this smoothing
problem) is
\[
X_t\sim
\mathcal N(x_0,\,2\kappa t).
\]

Suppose an exact later endpoint $X_T=b$ is learned.
This is an idealized observation. To define its conditional law despite
$\mathbb P(X_T=b)=0$, work on the Polish space
\[
\Wcal_T=\{w\in C([0,T],\mathbb R):w(0)=x_0\}
\]
with the uniform topology and its Borel sigma-field.

\paragraph{WREH dictionary for the endpoint experiment.}
The benchmark uses the following fixed objects. Here $\mu$ is the
unconditioned diffusion path law, and $0<\tau<T$ is a fixed earlier cut.
\begin{center}
\begin{tabularx}{\linewidth}{@{}lX@{}}
\toprule
Object & Declaration \\
\midrule
World / carrier & $w\in\Wcal_T$, a continuous path starting at $x_0$;
  $\mu$ supplies its prior probability law. \\
Protocol / response & Endpoint evaluation $R_{P_T}(w)=w(T)$. \\
Data & $D_b$ records the exact endpoint $b$; $x_0$ is fixed in the carrier. \\
Completion fibre & $\Ccal(D_b)=\{w\in\Wcal_T:w(T)=b\}$,
  also denoted $\Ccal(x_0,b)$. \\
Full pre-endpoint past & $\pi_-(w)=w|_{[0,T)}$,
  with codomain $\pi_-(\Wcal_T)$. \\
Fixed earlier fragment & $\pi_\tau(w)=w|_{[0,\tau]}$,
  with codomain $C_{x_0}([0,\tau],\mathbb R)$. \\
Conditional path law & $K_b$, a probability measure concentrated on
  $\Ccal(D_b)$; the fibre is a set, whereas the bridge is a law on paths. \\
\bottomrule
\end{tabularx}
\end{center}
The two historical projections answer different compatibility questions,
as explained in Remark~\ref{rem:bridge-compatibility}.

For every $b\in\mathbb R$, let $K_b$ be the law of
\[
U_t^b=x_0+\frac tT(b-x_0)
+\sqrt{2\kappa}\left(B_t-\frac tT B_T\right),
\qquad 0\le t\le T.
\]
The residual Gaussian bridge has zero covariance with $B_T$ at every
time. Thus its finite-dimensional vectors are independent of $B_T$;
continuity and evaluations at rational times extend this to the path
sigma-field. Since $X_t= x_0+t(X_T-x_0)/T+
\sqrt{2\kappa}(B_t-tB_T/T)$, it follows that
\[
\mathbb P(X\in F,\ X_T\in E)
=\int_E K_b(F)\,\mathbb P_{X_T}(db)
\]
for Borel $F\subseteq\Wcal_T$ and $E\subseteq\mathbb R$.
Translation by the continuous mean path makes $b\mapsto K_b$ weakly
continuous, in particular a measurable probability kernel.
An arbitrary regular conditional law is determined only
$\mathbb P_{X_T}$-almost everywhere; this explicit kernel selects a
version for every $b$. The notation $X\mid X_T=b$ refers to $K_b$,
not to division by a zero event probability. This residual construction
is classical \cite[Theorem 22.1 and the following remark]{Pitman2003}.
Its covariance is
\[
\operatorname{Cov}_{K_b}(X_s,X_t)
=2\kappa\left(\min(s,t)-\frac{st}{T}\right).
\]

\begin{theorem}[Past uncertainty contracts but does not collapse]
\label{thm:brownian-bridge}
Conditioned on the later endpoint \(X_T=b\),
\[
X_t\mid X_T=b
\sim
\mathcal N\!\left(
x_0+\frac{t}{T}(b-x_0),
\;
2\kappa\,\frac{t(T-t)}{T}
\right).
\]
Therefore
\[
0<
2\kappa\,\frac{t(T-t)}{T}
<
2\kappa t
\qquad
(0<t<T).
\]
The later exact observation strictly reduces uncertainty about the earlier state, but it does
not determine that earlier state uniquely.
\end{theorem}

\begin{proof}
The pair \((X_t,X_T)\) is jointly Gaussian with
\[
\operatorname{Var}(X_t)=2\kappa t,
\qquad
\operatorname{Var}(X_T)=2\kappa T,
\qquad
\operatorname{Cov}(X_t,X_T)=2\kappa t.
\]
Applying Theorem~\ref{thm:gaussian-contraction} gives
\[
\operatorname{Var}(X_t\mid X_T)
=
2\kappa t-\frac{(2\kappa t)^2}{2\kappa T}
=
2\kappa\,\frac{t(T-t)}{T}.
\]
The Gaussian conditional-mean formula gives the stated mean.
\end{proof}

\begin{corollary}[A full past path remains non-unique]
\label{cor:path-nonuniqueness}
The kernel $K_b$ is a Brownian-bridge law on continuous paths
and is not concentrated on a single trajectory. It assigns probability
zero to every countable family of trajectories.
\end{corollary}
\begin{proof}
Fix $t_*\in(0,T)$. Its evaluation marginal under $K_b$ is a
nondegenerate Gaussian law. Every singleton path imposes a singleton
evaluation value, which has probability zero. Countable unions of such
singleton paths also have probability zero.
\end{proof}

\begin{remark}
\label{rem:bridge-compatibility}
In this benchmark declare hard-constraint compatibility by
\[
\Ccal(x_0)=\Wcal_T,\qquad
\Ccal(x_0,b)=\{w\in\Wcal_T:w(T)=b\}.
\]
The endpoint fibre is a proper subset of $\Wcal_T$, and $K_b$ is
concentrated on it. These are compatibility classes, not sets of paths
with positive singleton probability.
For the full past before $T$, take $\pi_-(w)=w|_{[0,T)}$ and its
codomain to be $\pi_-(\Wcal_T)$. Continuity makes this restriction
injective. Its endpoint-conditioned image is therefore a proper
subset of the prior image, characterized by limit $b$ as $t\uparrow T$.
It remains uncountable by the preceding corollary.

By contrast, for any fixed $\tau<T$, the projected compatibility set
of fragments on $[0,\tau]$ is unchanged: every continuous fragment
starting at $x_0$ can be extended continuously to $b$ at $T$.
Its projected conditional law nevertheless has reduced variance at
each $0<t\le\tau$. Thus strict whole-history refinement does not imply
strict refinement of every earlier historical projection. No claim
about topological support equality is inferred from positive variance.
\end{remark}

\subsection{Noisy later observation}

Let the later measurement be
\[
Y_T=X_T+\eta,
\qquad
\eta\sim\mathcal N(0,R),
\qquad
R>0,
\]
independent of the diffusion.

\begin{proposition}[Noisy retrospective contraction]
\label{prop:noisy-contraction}
For \(0<t<T\),
\[
\operatorname{Var}(X_t\mid Y_T)
=
2\kappa t\left(1-\frac{2\kappa t}{2\kappa T+R}\right),
\]
and
\[
0<
\operatorname{Var}(X_t\mid Y_T)
<
2\kappa t.
\]
As \(R\downarrow0\), the formula converges to the Brownian-bridge
variance; as $R\to\infty$, it converges to the prior variance.
The continuous Gaussian kernel has mean
\[
\mathbb E[X_t\mid Y_T=y]
=x_0+\frac{2\kappa t}{2\kappa T+R}(y-x_0).
\]
\end{proposition}

\begin{proof}
Use Theorem~\ref{thm:gaussian-contraction} with
\[
\operatorname{Cov}(X_t,Y_T)=2\kappa t,
\qquad
\operatorname{Var}(Y_T)=2\kappa T+R.
\]
The inequalities follow because \(\kappa,t,R>0\) and \(t<T\).
\end{proof}

\begin{remark}[Noisy compatibility and support]
For observed $y$, the path likelihood is
\[
L_y(w)=\frac{1}{\sqrt{2\pi R}}
\exp\!\left[-\frac{(y-w(T))^2}{2R}\right]>0.
\]
It is bounded and has a finite positive prior integral. Bayes' formula
therefore defines a posterior equivalent to the prior path measure:
both have the same null sets and the same topological support.
If worlds include measurement noise, exact observation constrains
$(w,\eta)$ by $w(T)+\eta=y$, but its projection onto $w$ is all of
$\Wcal_T$, since $\eta=y-w(T)$ is always available.
\end{remark}

\begin{remark}[Physical reading]
A later noisy measurement improves retrodiction of an earlier diffusive state without making
the past exact. This is standard fixed-interval smoothing behavior, not a WREH discovery.
The WREH role of the benchmark is to make the status-of-the-past distinction operational:
later evidence can sharpen a past posterior while residual historical ambiguity remains.
\end{remark}

\section{Construction order versus internal time}
\begin{remark}[Inference order and temporal order]
The definitions supply an evidence-refinement order and, separately,
a declared temporal model. No map identifying the two orders is
specified. Any physical identification requires an additional stated
bridge. This is a scope convention, not a general impossibility theorem.
\end{remark}

\section{Prior-art boundary}
Inverse-limit history systems and dyadic solenoids are established
objects \cite[Section 1]{ClarkFokkink2004}. Gaussian conditioning
\cite[Appendix A.1]{SarkkaSvensson2023}, Brownian bridges
\cite{Pitman2003}, and retrospective smoothing
\cite{RauchTungStriebel1965} are classical tools.
The image, realizability-gate, and refinement propositions are
elementary consequences of the definitions. The structural-wall
transfer follows by conjugating local trivializations and inherits
the WR-III notion; it is not a new bifurcation theorem. This version organizes these results within
the WREH history-projection framework; it claims no theorem priority
for them. Originality of the wider classification has not been
certified by an exhaustive prior-art comparison.
No projective-extension theorem is used: existence of Brownian motion
is assumed in the benchmark, and its conditional kernel is constructed
explicitly.

\section{Open obligations}
\begin{enumerate}
\item complete the broader prior-art comparison with observability,
identifiability, natural extensions, retrodiction, and reconstruction
under partial observations; the source mapping for the specific
benchmarks is recorded in the accompanying repair audit;
\item retain the distinction between a framework exposition and a new
technical theorem; upstream DOI records have been checked, while a full
cross-manuscript duplication audit remains separate;
\item give a several-observation-time example in a declared linear
Gaussian state-space model, using the standard Rauch--Tung--Striebel
(RTS) backward recursion \cite[Section 3.2]{RauchTungStriebel1965}
\cite[Section 12.2, Theorem 12.2]{SarkkaSvensson2023}, and compare
exact-set versus probabilistic historical uncertainty. This is an
application obligation, not a new smoothing algorithm; nonlinear
models require separate treatment;
\item finite-resolution observable-past equivalence;
\item a non-duplicative topological or metric theory of past-completion sets;
\item an HTML demo of backward-history branching and evidence refinement.
\end{enumerate}

\section*{Conclusion}
Knowing a present response, knowing that a global completion exists, knowing a unique hidden past, and knowing a unique observable past are distinct statements. An unrealizable full profile has an empty historical class only relative to the declared model. Structural walls transfer to historical families under the stated topological bridge, and may otherwise disappear under projection. Under accumulated hard constraints the compatible completion classes are nested, possibly equal; no rewriting operation is required. Noninvertible deterministic dynamics provides a clean benchmark in which the future is unique while the full past remains highly nonunique. In the Brownian endpoint benchmark, later observations reduce
interior-state variance without determining a unique path. Exact
endpoint data restrict full histories before $T$, but need not exclude
any fragment on a fixed interval ending strictly before $T$. Noisy
Gaussian endpoint data preserve path-measure support. These conclusions
refer to different declared mathematical objects.

\paragraph{Review status.}
This is a revised review draft responding to Reviewer 1, not a
publication-ready release.
The broader prior-art and cross-manuscript checks listed above remain
open. Content licence: CC BY 4.0, subject to the repository's
third-party-material exclusions.

\begin{thebibliography}{99}
\bibitem{WRI} A. A. Malachevsky, \emph{Response-Defined World Spaces: Response Equivalence, Finite-Resolution Separation, and Identifiability of Global Realizations}, WR-I, WREH (2026). DOI: 10.5281/zenodo.23210258.
\bibitem{WRII} A. A. Malachevsky, \emph{Finite Consistency and Global World Realizability}, WR-II, WREH (2026). DOI: 10.5281/zenodo.23224154.
\bibitem{WRIII} A. A. Malachevsky, \emph{Geometry of Admissible World Fibres: Refinement, Rigidity and Realizability Walls}, WR-III, WREH (2026). DOI: 10.5281/zenodo.23228682.

\bibitem{RauchTungStriebel1965}
H. E. Rauch, F. Tung, and C. T. Striebel,
\emph{Maximum Likelihood Estimates of Linear Dynamic Systems},
AIAA Journal 3(8) (1965), 1445--1450.
DOI: 10.2514/3.3166.

\bibitem{SarkkaSvensson2023}
S. S\"arkk\"a and L. Svensson,
\emph{Bayesian Filtering and Smoothing}, 2nd ed.,
Cambridge University Press, 2023.
DOI: 10.1017/9781108917407.

\bibitem{KaratzasShreve1998}
I. Karatzas and S. E. Shreve,
\emph{Brownian Motion and Stochastic Calculus}, 2nd ed.,
Springer, 1998.
DOI: 10.1007/978-1-4612-0949-2.
\bibitem{ClarkFokkink2004}
A. Clark and R. Fokkink,
\emph{Embedding solenoids},
Fundamenta Mathematicae 181(2) (2004), 111--124.
DOI: \href{https://doi.org/10.4064/fm181-2-2}{10.4064/fm181-2-2}.

\bibitem{Pitman2003}
J. W. Pitman, \emph{Brownian Bridge},
Stat205B: Probability Theory, Lecture 22, Spring 2003,
University of California, Berkeley; scribe T. Chen.
\url{https://www.stat.berkeley.edu/~pitman/s205s03/b_bridge.pdf}.
\end{thebibliography}
\end{document}
